\documentclass[12pt]{article}
\usepackage{amsmath, amssymb, geometry, enumitem}
\geometry{margin=1in}
\setlength{\parindent}{0pt}
\setlength{\parskip}{0.5em}

\title{ECON 101: Macroeconomics \\ Quiz 4}
\author{Instructor: Muhebullah Karimzada}
\date{March 31,  2026}

\begin{document}
\maketitle







\begin{enumerate}[label=\textbf{\arabic*.}]

\item Which of the following is the best example of money serving as a \textbf{unit of account}?
\begin{enumerate}[label=(\Alph*)]
    \item A household keeps part of its wealth in cash.
    \item A consumer uses a debit card to buy groceries.
    \item A car dealership lists the price of a car as \$28{,}000.
    \item A worker exchanges labour services for wages.
\end{enumerate}

\item Suppose the desired reserve ratio is \(r_d = 0.20\). The simple deposit multiplier is:
\begin{enumerate}[label=(\Alph*)]
    \item \(0.20\)
    \item \(4\)
    \item \(5\)
    \item \(20\)
\end{enumerate}

\item If the Bank of Canada sells government securities, the most direct immediate effect is:
\begin{enumerate}[label=(\Alph*)]
    \item Bank reserves increase and the money supply rises.
    \item Bank reserves decrease and the money supply falls.
    \item The desired reserve ratio falls automatically.
    \item The inflation target rises.
\end{enumerate}

\item Which of the following would most likely be included in the narrow definition of the money supply?
\begin{enumerate}[label=(\Alph*)]
    \item A five-year government bond
    \item A savings account with withdrawal restrictions
    \item Currency held by the public
    \item Shares of a mutual fund
\end{enumerate}
\newpage
\item When commercial banks hold \textbf{excess reserves}, the actual increase in the money supply will typically be:
\begin{enumerate}[label=(\Alph*)]
    \item greater than the maximum predicted by the simple multiplier
    \item equal to the maximum predicted by the simple multiplier
    \item smaller than the maximum predicted by the simple multiplier
    \item unaffected by the reserve ratio
\end{enumerate}

\end{enumerate}

\hrule
\vspace{1em}

\section*{ Problem }



Assume the following:
\begin{itemize}
    \item A customer deposits \$8{,}000 in Bank A.
    \item The desired reserve ratio is \(r_d = 0.10\).
    \item Banks lend out all excess reserves.
    \item The public redeposits all loan proceeds into the banking system.
\end{itemize}

\begin{enumerate}
    \item[(a)] Calculate the amount of desired reserves that Bank A must hold from the initial deposit. (2 marks)
    \item[(b)] Calculate the amount Bank A can initially lend out. (2 marks)
    \item[(c)] Calculate the simple deposit multiplier. (2 marks)
    \item[(d)] Calculate the maximum total increase in deposits in the banking system resulting from the initial \$8{,}000 deposit. (3 marks)
    \item[(e)] Calculate the maximum total increase in loans in the banking system. (3 marks)
    \item[(f)] Suppose instead that banks decide to hold some excess reserves and households keep some cash rather than redepositing all funds. Explain how and why the actual increase in deposits would differ from your answer in part (d). (3 marks)
\end{enumerate}

\hrule
\vspace{1em}

\begin{center}
\textit{End of Quiz}
\end{center}

\end{document}